\chapter*{第 2 章 流体静力学}\addcontentsline{toc}{chapter}{第 2 章 流体静力学}\markboth{第 2 章 流体静力学}{}

\begin{tishi}
本章题目取自孔珑《工程流体力学（第四版）》第 3 章"流体静力学"课后习题（原书题号 3-1 至 3-31），与赵琴教材（精讲本）第 2 章"流体静力学"逐节对应：\S 2.1 流体静压强特性、\S 2.2 流体平衡微分方程、\S 2.3 重力场中流体静压强分布（含压强表示法与静压强的测量）、\S 2.4 流体的相对平衡、\S 2.5 液体作用在平面上的总压力、\S 2.6 液体作用在曲面上的总压力（含浮力）。31 道题全部落在西华 818 考纲（赵琴教材第 1 至 8 章）范围内，故\emph{全数收录}，无略去题。每题题首标注\textbf{孔珑原书题号}，题后给出重要度（星级）、依据与难度；原书（或其 OCR 文本）中明显不合物理量纲或不便阅读之处，在章末按题号列出说明。
\end{tishi}

\begin{ti}\sloppy\emergencystretch=2em
\item[\textbf{3-1}] 烟囱高 $H=20$ m，烟气温度 $t_g=300\,\mathrm{^\circ C}$，烟气密度按 $\rho_g=(1.25-0.0027t)$ kg/m$^3$ 计算（$t$ 为烟气温度，单位 $\mathrm{^\circ C}$），烟囱外空气密度 $\rho_a=1.29$ kg/m$^3$。求烟囱底部内外的压强差。
\tuyi{烟囱内为高温烟气柱、外侧为常温空气柱，二者分别视为静止气体；两气柱在烟囱顶部相通，故以顶部为压强相同的基准面，比较底部内外压强。}
\xstarfull{4}\quad\yiju{E6 教材 2.3 节重力场中流体静压强分布；E2 考情（第 1 至 3 章偏概念与基本公式运用）}\quad\nandu{易}

\item[\textbf{3-2}] 一直煤气管沿铅垂方向敷设，上、下两截面相距 $H=20$ m，两截面处各装一个 U 形管测压计，管内装水（$\rho_w=1000$ kg/m$^3$）；管外空气密度 $\rho_a=1.28$ kg/m$^3$，测得下截面测压管读数 $h_1=100$ mm、上截面读数 $h_2=115$ mm。忽略测压计中气柱的影响，求管内煤气的密度 $\rho_g$。
\tuyi{两个 U 形管分别接在煤气管的上、下截面，两读数之差反映 20 m 气柱自重与管内煤气柱自重之差。}
\xstarfull{4}\quad\yiju{E6 教材 2.3 节重力场中流体静压强分布（U 形管测压计）；E2 考情（基本公式运用）}\quad\nandu{中}

\item[\textbf{3-3}] 充满水的 $A$、$B$ 两容器用 U 形压差计相连，压差计内装水银（$\rho_{Hg}=13600$ kg/m$^3$），测得两水银面高差 $h=15$ cm。求 $A$、$B$ 两容器的压强差。
\xstarfull{4}\quad\yiju{E6 教材 2.3 节重力场中流体静压强分布（压差计）；E4 2015 年填空题（U 形水银压差计）}\quad\nandu{易}

\item[\textbf{3-4}] 容器 $A$ 与一 U 形管测压计相接，测压计内装水银（$\rho_{Hg}=13600$ kg/m$^3$）。已知各液面高程（自同一基准面量起）为 $h_1=0.25$ m、$h_2=1.61$ m、$h_3=1$ m，$\rho_w=1000$ kg/m$^3$，大气压 $p_a=101325$ Pa。求容器内水面的绝对压强与真空度。
\tuyi{$h_1$ 为容器 $A$ 内水面的高程，$h_2$、$h_3$ 为测压计 U 形管内两水银面的高程（$h_2>h_3$，$h_2$ 一侧管内的水柱与容器 $A$ 相通，$h_3$ 一侧管口通大气）。}
\xstarfull{4}\quad\yiju{E6 教材 2.3 节（测压管与压差计）；E2 考情（基本公式运用）}\quad\nandu{中}

\item[\textbf{3-5}] 直径 $d=0.4$ m 的圆柱形容器，盖上加作用力 $F=5788$ N；容器内盛有 $h_1=30$ cm 的油（$\rho_{oil}=800$ kg/m$^3$）和 $h_2=50$ cm 的水，容器底部用一 U 形水银测压计测压（$\rho_{Hg}=13600$ kg/m$^3$）。求测压计中水银柱的高差 $H$。
\xstarfull{4}\quad\yiju{E6 教材 2.3 节（重力场中静压强分布与等压面）；E2 考情（基本公式运用）}\quad\nandu{中}

\item[\textbf{3-6}] 密封盛水容器接有两根 U 形水银测压管（$\rho_{Hg}=13600$ kg/m$^3$）。上管水银面距容器自由液面 $h_1=60$ cm，管内水银柱高 $h_2=25$ cm；下管内水银柱高 $h_3=30$ cm。求下管水银面距容器自由液面的高度 $h_4$。
\tuyi{两根测压管的上端均通大气，容器液面上方为密闭气体（压强 $p_0$）。}
\xstarfull{4}\quad\yiju{E6 教材 2.3 节（多支测压管）；E2 考情（基本公式运用）}\quad\nandu{中}

\item[\textbf{3-7}] 封闭容器内盛有油和水，油层厚 $h_1=30$ cm（$\rho_{oil}=800$ kg/m$^3$）；容器侧壁下部的 U 形水银测压管（$\rho_{Hg}=13600$ kg/m$^3$）中，水银面低于容器内油水分界面 $h_2=50$ cm、低于油面 $h_3=40$ cm，测压管另一端通大气（$p_a=101325$ Pa）。求油面上的计示压强及绝对压强。
\xstarfull{4}\quad\yiju{E6 教材 2.3 节；E2 考情（第 1 至 3 章偏概念与基本公式）}\quad\nandu{易}

\item[\textbf{3-8}] 水压机的大活塞上作用力 $F_1=4905$ N，作用在杠杆上的力 $F_2=147$ N，杠杆的短臂 $a=15$ cm、长臂 $b=75$ cm，小活塞直径 $d_1=5$ cm，考虑摩擦的力的校正系数 $\eta=0.9$。求大活塞的直径 $d_2$。
\xstarfull{4}\quad\yiju{E6 教材 2.1 节（帕斯卡定律与静压强的等值传递）；E2 考情（基本公式运用）}\quad\nandu{中}

\item[\textbf{3-9}] 双液式微压计的两杯直径 $d_1=50$ mm，U 形管直径 $d_2=5$ mm，$A$ 杯内装酒精（$\rho_1=870$ kg/m$^3$），$B$ 杯内装煤油（$\rho_2=830$ kg/m$^3$），测得两液体的分界面上升 $h=280$ mm。求 $A$、$B$ 两杯液面上的压强差。
\tuyi{两杯液面在测压过程中分别下降、上升同样大小的 $\delta$（由体积守恒确定），分界面上升 $h$。}
\xstarfull{4}\quad\yiju{E6 教材 2.3 节（双液微压计）；E2 考情（基本公式运用）}\quad\nandu{难}

\item[\textbf{3-10}] 用复式水银测压计测量锅炉水面上的蒸汽压强，各液面高程（自同一基准面量起）为 $H=3$ m、$h_1=1.4$ m、$h_2=2.5$ m、$h_3=1.2$ m、$h_4=2.3$ m，$\rho_{Hg}=13600$ kg/m$^3$，$\rho_w=1000$ kg/m$^3$，大气压 $p_a=101325$ Pa。求锅炉水面上的蒸汽绝对压强。
\tuyi{$H$ 为锅炉水面高程；测压计由两段水银柱与两段水柱串联而成，$h_1$ 至 $h_4$ 为各水银面高程。}
\xstarfull{4}\quad\yiju{E6 教材 2.3 节（复式测压计）；E2 考情（基本公式运用）}\quad\nandu{中}

\item[\textbf{3-11}] 倾斜微压计内装酒精（$\rho=810$ kg/m$^3$），倾斜角 $\alpha=45^\circ$，读数 $l=20$ cm。求所测压强差；若改用 $\rho_w=998$ kg/m$^3$ 的水、倾斜角改为 $30^\circ$，读数 $l'$ 为多少？
\xstarfull{4}\quad\yiju{E6 教材 2.3 节（倾斜微压计）；E2 考情（基本公式运用）}\quad\nandu{易}

\item[\textbf{3-12}] 直线行驶的汽车上放置一内装液体的 U 形管，两支管沿行驶方向相距 $l=500$ mm。汽车以加速度 $a=0.5$ m/s$^2$ 沿水平直线行驶，求两支管中的液面高度差。
\xstarfull{3}\quad\yiju{E6 教材 2.4 节流体的相对平衡（等加速直线运动）；E2 考情（第 1 至 3 章偏概念）}\quad\nandu{易}

\item[\textbf{3-13}] 油罐车内装着密度 $\rho=1000$ kg/m$^3$ 的液体，以水平速度 $u=36$ km/h 行驶，油罐车尺寸为 $D=2$ m、$h=0.3$ m、$l=4$ m。车在某一时刻开始减速，经 $100$ m 距离后完全停下。若为均匀制动，求作用在侧面 $A$（前端圆面）上的力。
\tuyi{油罐为卧式圆柱罐（直径 $D$、长 $l$）；$h$ 为自由液面高出罐轴线的高度（原书解答中液面高出罐轴线 $D/2+h=1.3$ m），制动时液面因惯性向前倾斜，取前端圆面为受力的积分面。}
\xstarfull{4}\quad\yiju{E6 教材 2.4 节（等加速直线运动容器）；E2 考情（基本公式运用）}\quad\nandu{中}

\item[\textbf{3-14}] 一正方形容器放在小车上，容器底边 $b=200$ mm、质量 $m_1=4$ kg，内盛水深 $h=150$ mm；小车质量 $m_3=6$ kg，用质量 $m_2=25$ kg 的重物经滑轮牵引小车沿水平面滑动，容器底与平面间的摩擦系数 $\mu=0.3$。求容器中水刚好不溢出时容器的最小高度 $H$。
\tuyi{重物经定滑轮牵引小车，容器内水面因水平加速度而倾斜；水的体积不变，前后液面高差绕容器中点对称。}
\xstarfull{3}\quad\yiju{E6 教材 2.4 节（等加速直线运动容器）；E2 考情（第 1 至 3 章偏概念与基本公式）}\quad\nandu{中}

\item[\textbf{3-15}] 一盛水容器以加速度 $a=4.9$ m/s$^2$ 铅垂向下作等加速运动，运动前水深 $h=2$ m（$\rho=1000$ kg/m$^3$，大气压 $p_a=101325$ Pa）。(1) 求容器底部的流体绝对静压强；(2) 加速度为何值时容器底部所受压强为大气压强；(3) 加速度为何值时容器底部的绝对静压强等于零（提示：本题铅垂方向单位质量力 $f_z=a-g$）。
\xstarfull{3}\quad\yiju{E6 教材 2.4 节（容器作等加速运动）；E2 考情（第 1 至 3 章偏概念与基本公式）}\quad\nandu{易}

\item[\textbf{3-16}] 一圆柱形容器直径 $d=300$ mm、高 $H=500$ mm，内盛水深 $h=300$ mm，绕铅垂轴旋转。(1) 求水刚好不溢出容器时的转速 $n$；(2) 求水刚好露出底面时的转速 $n$，并求停止旋转后的水深。
\tuyi{旋转时自由液面为抛物面，抛物面上、下顶点与静止液面之间的体积关系由容器内水量不变确定。}
\xstarfull{3}\quad\yiju{E6 教材 2.4 节（等角速度旋转容器）；E2 考情（第 1 至 3 章偏概念与基本公式）}\quad\nandu{中}

\item[\textbf{3-17}] 用离心浇铸机铸造车轮，铁水密度 $\rho=7138$ kg/m$^3$，其他尺寸 $h=200$ mm、$d=900$ mm，上箱及砂重 $G=10$ kN，转速 $n=600$ r/min。浇铸前铁水自由面为水平面且距上箱底面 $h$，求：(1) 圆筒壁面上铁水的计示压强；(2) 螺栓群所受的总拉力。
\xstarfull{4}\quad\yiju{E6 教材 2.4 节（等角速度旋转容器及其工程应用）；E2 考情（基本公式运用）}\quad\nandu{难}

\item[\textbf{3-18}] 一圆柱形容器直径 $d=1.2$ m，充满水并绕铅垂轴等角速度旋转；在盖顶上 $r_0=0.43$ m 处安装一开口测压管，管中水位高出盖顶 $h=0.5$ m。问容器转速 $n$ 为多少时顶盖所受的静水总压力为零？
\tuyi{测压管水位给出 $r=r_0$ 处的计示压强 $p_e=\rho g h$；顶盖上取半径为 $r$、宽 $\dd r$ 的圆环积分，令总压力为零求 $\omega$。}
\xstarfull{3}\quad\yiju{E6 教材 2.4 节（等角速度旋转容器）；E2 考情（基本公式运用）}\quad\nandu{难}

\item[\textbf{3-19}] 一圆柱形容器直径 $d=600$ mm、高 $H=500$ mm，内盛水至 $h=400$ mm，余下容积盛满密度 $\rho=800$ kg/m$^3$ 的油，容器顶盖中心有一小孔与大气相通。若容器绕其主轴旋转，问转速 $n$ 多大时油面开始接触到底板？并求此时顶盖和底板上的最大、最小计示压强（提示：油水分界面为等压面）。
\tuyi{旋转时油水分界面为抛物面：在轴心处触及底板、在 $r=r_1$ 处触及顶盖；由油体积不变确定 $r_1$。}
\xstarfull{3}\quad\yiju{E6 教材 2.4 节（等角速度旋转容器）；E2 考情（基本公式运用）}\quad\nandu{难}

\item[\textbf{3-20}] 斜壁上装一圆形闸门，直径 $d=0.5$ m，闸门顶端沿斜壁到水面的距离 $a=1$ m，斜壁与水平面成 $\alpha=60^\circ$。求水作用在闸门上的总压力及压力中心（压力中心以沿斜壁距水面的距离表示）。
\xstarfull{5}\quad\yiju{E6 教材 2.5 节液体作用在平面上的总压力（$P=\rho g h_C A$、压力中心 $y_D=y_C+J_C/(y_C A)$）；E1 考情（流体静力学为考试计算题基础入门、重点掌握方法）；E4 2015 年计算题（闸门静水总压力）}\quad\nandu{易}

\item[\textbf{3-21}] 一绕铰链 $O$ 转动的自动开启式水闸，闸门与水平面成 $\alpha=60^\circ$，闸门一侧水深 $H=2$ m、另一侧水深 $h=0.4$ m。求闸门开始自动开启（两侧水压力对 $O$ 的力矩相等）时铰链 $O$ 至闸门下端的距离 $x$。
\tuyi{闸门绕 $O$ 铰转动，$O$ 位于斜壁面上；两侧水压力均垂直于闸门面，故它们对 $O$ 的力臂即沿闸门面量得的距离，压力中心距各自液面分别为 $2H/(3\sin\alpha)$ 与 $2h/(3\sin\alpha)$。}
\xstarfull{5}\quad\yiju{E6 教材 2.5 节（平面总压力与压力中心、力矩平衡）；E1 考情（计算题基础入门）；E2 考情（计算题占 50\% 分值）}\quad\nandu{中}

\item[\textbf{3-22}] 水作用在如图所示的 $3/4$ 圆柱面 $ABCD$ 上，试就 (a)、(b)、(c) 三种开口测压管的液面位置 $V_1$、$V_2$、$V_3$，画出压力体及所对应的总压力铅垂分力的方向。
\tuyi{圆柱面 $ABCD$ 为 $3/4$ 圆弧；三种情形下测压管液面分别为 $V_1$、$V_2$、$V_3$，压力体为受力曲面、自由液面（或其延长面）与各点铅垂母线所围体积。}
\xstarfull{5}\quad\yiju{E6 教材 2.6 节压力体虚实与方向；E1 考情（流体静力学为重点、压力体为常考画图内容）}\quad\nandu{中}

\item[\textbf{3-23}] 一盛水的球形容器直径 $d=2$ m，下半个球固定不动，上半个球用螺栓与下半个球连接，球内充满水。不求及结构自重，求作用于螺栓上的力。
\tuyi{对上半个球内的水体列铅垂方向平衡；两半球在水平方向的水压力相互抵消。}
\xstarfull{5}\quad\yiju{E6 教材 2.6 节液体作用在曲面上的总压力（铅垂分力与压力体）；E1 考情（计算题基础入门）}\quad\nandu{中}

\item[\textbf{3-24}] 一储水设备在 $C$ 点测得绝对压强 $p_C=196120$ Pa，$h=2$ m，$R=1$ m，求作用于半球 $AB$ 的总压力（$\rho_w=1000$ kg/m$^3$，大气压 $p_a=101325$ Pa）。
\tuyi{$C$ 点为设备内水面上的测压点；$h$ 为半球底面（$AB$ 所在平面）的淹没深度，$R$ 为半球半径。}
\xstarfull{5}\quad\yiju{E6 教材 2.6 节（曲面总压力的水平分力与铅垂分力）；E1 考情（计算题基础入门、重点掌握方法）}\quad\nandu{难}

\item[\textbf{3-25}] 一扇形闸门（圆弧面闸门），宽度 $B=1$ m，圆心角 $\alpha=45^\circ$，上游水头 $H=3$ m。求水对闸门作用力的大小及方向。
\tuyi{圆弧面上各点的压强沿半径指向圆心，故合力作用线通过圆心 $O$；压力体由圆弧面、自由液面及过圆弧下端的铅垂线围成。}
\xstarfull{5}\quad\yiju{E6 教材 2.6 节（曲面总压力：水平分力与铅垂分力）；E5 题库计算题（弧形闸门求水平分力、铅垂分力与力矩）；E1 考情（计算题基础入门）}\quad\nandu{难}

\item[\textbf{3-26}] 一扇形闸门，半径 $R=7.5$ m，挡着深度 $h=4.8$ m 的水，其圆心角 $\alpha=43^\circ$，旋转轴距渠底 $H=5.8$ m，闸门的水平投影 $CB=a=2.7$ m，闸门宽度 $B=6.4$ m。求作用在闸门上的总压力的大小和压力中心。
\tuyi{圆弧面对应的圆心 $O$ 在闸门上方；压力体由圆弧面、自由液面及过圆弧两端的铅垂母线围成；压力中心为过 $O$ 的合力作用线与圆弧面的交点。}
\xstarfull{5}\quad\yiju{E6 教材 2.6 节（曲面总压力与压力中心）；E5 题库计算题（弧形闸门求水平分力、铅垂分力与力矩）}\quad\nandu{难}

\item[\textbf{3-27}] 盛有水的容器底部有圆孔口，用空心金属球体封闭。该球体的重力 $W=2.452$ N、半径 $r=4$ cm，孔口直径 $d=5$ cm，水深 $H=20$ cm。求提起该球体所需之最小力 $F$。
\tuyi{球体座在孔口上：球面与孔口圆周相切，球心在孔口平面上方 $r-h$ 处（$h$ 为球缺高）；孔口平面以上的球面受水压力向下，球体伸入孔内的球缺部分受水压力向上，二者均用压力体计算。}
\xstarfull{5}\quad\yiju{E6 教材 2.6 节（曲面总压力与压力体、球缺体积）；E1 考情（计算题基础入门）}\quad\nandu{难}

\item[\textbf{3-28}] 汽油箱底部有锥阀，其尺寸为 $d_1=100$ mm、$d_2=50$ mm、$d_3=25$ mm、$a=100$ mm、$b=50$ mm，汽油密度 $\rho=830$ kg/m$^3$，略去阀芯自重和运动时的摩擦阻力。试确定：(1) 当压强表读数为 $9.806\times10^3$ Pa 时，提升阀芯所需的初始力 $F$；(2) $F=0$ 时箱中空气的计示压强 $p_e$ 等于多少。
\tuyi{锥阀上、下两侧分别对应两个压力体；上侧压力体为环形柱体（外径 $d_1$、内径 $d_3$，高为自由液面至阀尖的距离），下侧压力体由两段锥体与环形柱体组合而成。}
\xstarfull{5}\quad\yiju{E6 教材 2.6 节（曲面总压力与压力体）；E2 考情（计算题占 50\% 分值）}\quad\nandu{难}

\item[\textbf{3-29}] 一直径 $d=1$ m、高 $H=1.5$ m 的圆柱形容器内充满密度 $\rho=900$ kg/m$^3$ 的液体，顶盖中心开孔通大气。容器绕中心轴以 $n=50$ r/min 的转速旋转，求容器的上盖、底面和侧面所受的液体总压力。
\xstarfull{4}\quad\yiju{E6 教材 2.4 节（等角速度旋转容器的压强分布）与 2.5 节、2.6 节（平面与曲面上的总压力）；E2 考情（基本公式运用）}\quad\nandu{难}

\item[\textbf{3-30}] 一转动桥梁支承于直径 $d_1=3.4$ m 的圆形浮筒上，浮筒漂浮于直径 $d_2=3.6$ m 的室内。试确定：(1) 无外载荷而只有桥梁和浮筒自身重力 $W=29.43\times10^4$ N 时，浮筒沉没在水中的深度 $H$；(2) 当桥梁上有外载荷 $F=9.81\times10^4$ N 时，浮筒增加的沉没深度 $h$。
\tuyi{浮筒在室内水中漂浮：浮力等于总重力；浮筒下降 $h$ 的同时室内水面上升 $\Delta h$，二者由室内水的体积不变联系。}
\xstarfull{5}\quad\yiju{E6 教材 2.6 节（浮力与阿基米德原理）；E1 考情（计算题基础入门、重点掌握方法）}\quad\nandu{中}

\item[\textbf{3-31}] 一钢筋混凝土沉箱，长 6 m、宽 5 m、高 5 m、底厚 0.5 m、侧壁厚 0.3 m，钢筋混凝土密度 $\rho_1=2400$ kg/m$^3$，海水密度 $\rho_2=1025$ kg/m$^3$，沉箱在海面上漂浮是否稳定？
\tuyi{先求沉箱的浮心与重心高度，再由稳心半径与初稳性高判断漂浮稳定性。}
\xstarfull{4}\quad\yiju{E6 教材 2.6 节（浮力、浮心与漂浮稳定性）；E2 考情（第 1 至 3 章偏概念与基本公式）}\quad\nandu{难}
\end{ti}

\begin{tishi}
\textbf{第 3-1 题（星级说明）}：本题是静压强基本方程 $p=p_0+\rho gh$ 在气体中的应用（烟囱自流通风），精讲第 2 章与之对应的考点"流体静力学基本方程与测压管水头"按 $\star$5 计；因本题属该考点的边角应用（不涉及测压管水头线），本册按 $\star$4 计。
\end{tishi}

\begin{tishi}
\textbf{第 3-12 题（数值说明）}：原书解答给出 $\Delta z=0.0225$ m，而按等压面方程 $\Delta z=(a/g)l=0.5\times0.5/9.8=0.0255$ m 应为 25.5 mm；原书此处疑为排印或识别之误，本册取 $25.5$ mm。
\end{tishi}

\begin{tishi}
\textbf{第 3-13 题（坐标口径说明）}：原书解答取 $z$ 轴向上、原点在罐轴线，并取 $x=0$ 处自由液面高出罐轴线 $D/2+h=1.3$ m、$p=p_a$，由此定积分常数 $C=p_a+\rho g\times1.3=114065$ Pa；罐前端面圆心处 $x=l=4$ m、$z=0$，得 $p_{c}=116065$ Pa。题给 $h=0.3$ m 即按"液面高出罐轴线 $D/2+h$"读取。
\end{tishi}

\begin{tishi}
\textbf{第 3-24 题（压力体说明）}：原书解答中压力体体积的表达式 OCR 残缺，可辨认出折算水柱高 $h_c=(p_C-p_a)/(\rho g)=9.672$ m、结果 $P_z=\rho gV_p=246.33$ kN、$P_x=0$。反推 $V_p=246330/9800\approx25.14$ m$^3$，与 $V_p=\pi R^2(h_c-h)+\pi R^3/3$ 相符，本册即按此式计算。
\end{tishi}

\begin{tishi}
\textbf{第 3-27 题（压力体说明）}：原书压力体公式的 OCR 文本（page-0039.txt、page-0040.txt）残缺，可辨认出球缺高 $h=r-\sqrt{r^2-(d/2)^2}=0.8775$ cm、$V_1=8.964$ cm$^3$、$V_2=127.46$ cm$^3$，以及算式 $F=W+\rho gV_2-\rho gV_1=3.613$ N。本册按原书数值与算式给出结果。
\end{tishi}

\begin{tishi}
\textbf{第 3-28 题（密度口径说明）}：原书解答在把压强表读数折算为液柱高时取 $\rho=1000$ kg/m$^3$（得 $9.806\times10^3/(1000\times9.8)=1.0$ m，再加 $b$ 得自由液面高 $1.05$ m），与题给汽油密度 $830$ kg/m$^3$ 不一致。本册按原书口径给出 $F=12.579$ N、$p_e=1251.46$ Pa；若严格按 $\rho=830$ kg/m$^3$ 计算，各压力体的液柱高与 $\rho g$ 均需相应换算，结果将明显不同（提示读者按题给密度复核）。
\end{tishi}

\begin{tishi}
\textbf{第 3-31 题（星级说明）}：漂浮稳定性（浮心、稳心、初稳性高）在精讲第 2 章中归入"浮力与阿基米德原理"，该考点按 $\star$5 计；本题只考稳定性判断、不含复杂积分，故本册按 $\star$4 计。
\end{tishi}

\begin{tishi}
\textbf{本章收录与略去情况}：孔珑第 3 章课后习题 3-1 至 3-31 共 31 题，考查内容为静压强特性、平衡微分方程、重力场压强分布、静压强测量、相对平衡、平面与曲面总压力、浮力与漂浮稳定性，全部与赵琴教材第 2 章对应，\emph{没有}超出西华 818 考纲（气体动力学、流体机械、明渠堰流等）的题目，故\emph{本章无略去的题号}。
\end{tishi}

\begin{tishi}
\textbf{本章 OCR 订正清单}（供校对原书时参考）：
\begin{xiaowen}
\item 3-1：烟气密度单位 OCR 作 kg/m$^2$，应为 kg/m$^3$；"烟肉"应为"烟囱"。
\item 3-8：题面与解答中密度、直径单位多处 OCR 作 kg/m$^2$、m$^2$，应分别为 kg/m$^3$、m。
\item 3-12：原书结果 0.0225 m，按 $\Delta z=(a/g)l$ 应为 0.0255 m（见上文说明）。
\item 3-17：数字 $2864.29$ kPa 与 $905.086$ kN 可辨，OCR 中 $\rho$ 的单位误作 kg/m$^2$。
\item 3-19：油水分界面触及底板的半径 OCR 作 $(2r)^2=0.144$，实为 $(2r_1)^2=0.144$ m$^2$，即 $r_1=0.1897$ m。
\item 3-25、3-26：压力体体积的积分表达式 OCR 残缺，数字结果（$11398.397$ N、$45549.24$ N、$53.732$ m$^3$、$526573.6$ N、$894055$ N）可辨。
\item 3-27：$V_1$ 的 OCR 文本作"$=$"处应为"$\approx$"，单位 cm$^3$ 偶作 cm$^2$；球缺体积公式为 $\pi h^2(r-h/3)$。
\item 3-28：汽油密度 OCR 作 kg/m$^2$，应为 kg/m$^3$；答案中 $V_{p1}$、$V_{p2}$ 的单位由 m$^3$ 误作 m$^2$。
\item 3-31：海水密度 OCR 作 kg/m，应为 kg/m$^3$。
\end{xiaowen}
\end{tishi}
